\documentclass[11pt]{article}

\usepackage[margin=1in]{geometry}
\usepackage{longtable}
\usepackage{array}
\usepackage{hyperref}

\title{RelaLeap Experiment Summary and Open Research Questions}
\author{Generated from the RelaLeap automation evidence trail}
\date{2026-06-23}

\begin{document}
\maketitle

\begin{abstract}
This report summarizes the RelaLeap experiments completed so far, with an
emphasis on what has actually been learned and what remains uncertain. The
goal is to give external LLMs and human reviewers enough context to suggest
routes toward more sizeable improvements. The current strongest empirical
findings are: (i) temporal-consistency clipped HEP is a robust, deployable
label-free support-stress mitigation, but its effect size is small; (ii)
simple residual-objective variants mostly fail to beat supervised cross-entropy
under the present gates; and (iii) the newest local diagnostic suggests that
support width, rather than column count or another CE-adjacent loss, may be the
current residual-layer bottleneck.
\end{abstract}

\section{Context and Vocabulary}

RelaLeap is an experimental residual-layer learning program. A frozen base
model is augmented with residual columns at one or more insertion sites. The
experiments ask whether residual columns can adapt the frozen model in useful
ways while preserving guardrails on ordinary behavior and artifact integrity.

The main recurring concepts are:

\begin{itemize}
  \item \textbf{Residual columns}: learned residual atoms/columns inserted into
  the frozen base computation.
  \item \textbf{Support}: the selected residual columns/atoms active for a
  given example. Most early runs used sparse support with \texttt{top\_k = 1}.
  \item \textbf{HEP}: hidden-state error propagation / settling at inference
  time, implemented as an $\alpha$ sweep over a local hidden-state update.
  \item \textbf{Support stress}: a diagnostic condition that intentionally
  creates nonzero support repicking so that support instability can be measured.
  \item \textbf{Pinned support}: HEP is run while holding the original support
  fixed. This is diagnostically useful, but early evidence showed it can harm
  loss when support really wants to move.
  \item \textbf{Clipped HEP}: HEP updates are clipped to keep logit movement and
  pinned-vs-repicked divergence inside explicit budgets.
  \item \textbf{Temporal-consistency HEP}: a deployable label-free HEP
  objective that uses consistency over temporal/sequence perturbations rather
  than supervised labels during settling.
  \item \textbf{Artifact gate}: each comparison must produce expected
  \texttt{summary.json}, \texttt{metrics.csv}, and \texttt{notes.md} files, and
  model/summary invariants must pass.
\end{itemize}

\section{Evidence Protocol}

The project now has a command-driven experiment harness, local and Colab
comparison paths, artifact-contract checks, and decision reports. This matters
because many observed effects are small, and because the Colab path is driven
through a real-Chrome bridge rather than a clean official API.

Important guardrails include:

\begin{itemize}
  \item Local and Colab comparisons should match when both are run.
  \item HEP $\alpha = 0$ must reproduce ordinary inference.
  \item Frozen base parameters must remain unchanged during residual training.
  \item Residual parameters must update during residual training.
  \item Nonzero HEP alphas are accepted only if they improve the relevant loss
  while staying inside ordinary-logit and support-stability budgets.
  \item Promotion decisions are made by explicit reports, not by ad hoc visual
  inspection.
\end{itemize}

\section{High-Level Timeline}

\begin{enumerate}
  \item Phase 0 established the local and Colab smoke harness.
  \item Pinned-support and support-stress diagnostics exposed support
  instability.
  \item Clipped HEP stabilized support-stress settling but did not by itself
  improve loss.
  \item A guided clipped HEP oracle showed that a helpful hidden-state signal
  exists, but it used supervised labels and therefore was not deployable.
  \item Label-free HEP objectives were tested. Prediction-entropy was stable
  but poor; temporal consistency was stable and consistently slightly helpful.
  \item Temporal-consistency clipped HEP passed multi-seed, validation,
  extended, larger-char, and tokenized local/Colab gates, and was promoted as
  the default support-stress mitigation.
  \item With temporal clipped HEP fixed, residual-objective variants were
  tested. PC-style, anchored-PC, confidence penalty, margin penalty, label
  smoothing, focal CE, and train-time temporal consistency did not justify a
  default change.
  \item The newest diagnostic fixed supervised CE and temporal clipped HEP,
  varied only column count and support width, and found that widened support
  gave the largest local gain observed so far in this phase.
\end{enumerate}

\section{Phase 0: Harness and Baseline HEP}

Phase 0 validated the basic experiment harness on Tiny Shakespeare character
runs. The baseline configurations were:

\begin{itemize}
  \item \texttt{char\_smoke}: supervised CE residual update.
  \item \texttt{char\_smoke\_pc}: PC-style logit-MSE residual update.
  \item \texttt{char\_smoke\_hep}: supervised CE residual update plus HEP
  $\alpha$ sweep.
\end{itemize}

The Phase 0 baseline used hidden dimension 32, sequence length 32, one
insertion site, 8 residual columns, 4 atoms per column, \texttt{top\_k = 1},
seed 1, and 10 residual training steps.

\begin{center}
\begin{tabular}{l l r r r}
\hline
Experiment & Objective & Initial loss & Final loss & Delta \\
\hline
\texttt{char\_smoke} & supervised CE & 3.61089253 & 3.56317067 & -0.04772186 \\
\texttt{char\_smoke\_pc} & PC logit MSE & 0.02878798 & 0.02869168 & -0.00009630 \\
\texttt{char\_smoke\_hep} & supervised CE & 3.61089253 & 3.56317067 & -0.04772186 \\
\hline
\end{tabular}
\end{center}

The HEP alpha sweep produced a meaningful initial result: $\alpha = 0.25$ was
accepted because it improved loss while staying within the ordinary-logit
delta budget. Larger alphas improved raw loss more but exceeded the budget.

\begin{center}
\begin{tabular}{r r r l}
\hline
$\alpha$ & Loss & Max logit delta & Gate result \\
\hline
0.00 & 3.56317067 & 0.00000000 & baseline \\
0.25 & 3.55195642 & 0.05185765 & accepted \\
0.50 & 3.54104090 & 0.10371491 & rejected \\
1.00 & 3.52012682 & 0.20742965 & rejected \\
\hline
\end{tabular}
\end{center}

\textbf{Interpretation.} HEP can improve the supervised CE loss, but the
unconstrained update is too aggressive. This motivated the later clipping and
support-stability work.

\section{Support Stress, Pinned Support, and Clipping}

Support-stress runs intentionally induced support repicking. Pinned support was
tested as a possible stabilizer. It exposed the problem clearly, but it did not
improve performance under the acceptance gate.

\begin{center}
\begin{tabular}{r r r r}
\hline
$\alpha$ & Loss & Improvement vs $\alpha=0$ & Pinned-vs-repicked delta \\
\hline
0.00 & 4.32304668 & 0.00000000 & 0.00000000 \\
0.25 & 4.71638441 & -0.39333773 & 0.83149838 \\
0.50 & 5.17755079 & -0.85450411 & 1.66299534 \\
1.00 & 6.23724604 & -1.91419935 & 3.32599020 \\
\hline
\end{tabular}
\end{center}

The pinned-support decision was \texttt{keep\_opt\_in}. The max support change
fraction was 0.78906250, and the max pinned-vs-repicked logit delta was
3.32599020. This is a large instability signal.

Clipped HEP reduced the pinned-vs-repicked logit delta dramatically, from
3.32599020 to 0.00235271, but simple clipping alone did not improve the
support-stress loss. The clipped-HEP decision was also \texttt{keep\_opt\_in}.

\textbf{Interpretation.} Support instability is real. Clipping can make the
settling dynamics safe, but a useful settling direction is still needed.

\section{Guided HEP Oracle and Label-Free HEP Signals}

The guided clipped HEP oracle used supervised information during settling. It
showed that a helpful signal exists:

\begin{center}
\begin{tabular}{r r r r}
\hline
$\alpha$ & Loss & Improvement vs $\alpha=0$ & Logit delta \\
\hline
0.00 & 4.32304668 & 0.00000000 & 0.00000000 \\
0.25 & 4.32230139 & 0.00074530 & 0.00084373 \\
0.50 & 4.32155704 & 0.00148964 & 0.00168708 \\
1.00 & 4.32006598 & 0.00298071 & 0.00337455 \\
\hline
\end{tabular}
\end{center}

This result was important but not directly deployable, because the oracle used
supervised labels. Two label-free alternatives were then tested:

\begin{itemize}
  \item Prediction-entropy clipped HEP: deployable and stable, but worsened
  supervised support-stress loss in the local smoke comparison.
  \item Temporal-consistency clipped HEP: deployable, stable, and consistently
  slightly helpful under support stress.
\end{itemize}

\section{Temporal-Consistency Clipped HEP}

Temporal-consistency clipped HEP became the strongest early mechanism result.
It was selected across seed-smoke, validation, extended, larger-character, and
tokenized local/Colab evidence.

\subsection{Multi-Seed Smoke Evidence}

The multi-seed aggregate included 8 reports: seeds 1--4, local and Colab. All
selected temporal consistency. The mean loss improvement from $\alpha=0$ was
0.00016832. The max ordinary-logit delta was 0.00176525, and the max
pinned-vs-repicked delta was 0.00254941.

\begin{center}
\begin{tabular}{r l r r r}
\hline
Seed & Backend & Accepted $\alpha$ & Loss improvement & Logit delta \\
\hline
1 & local & 1.00 & 0.00020361 & 0.00176525 \\
1 & Colab & 1.00 & 0.00020313 & 0.00176525 \\
2 & local & 1.00 & 0.00012970 & 0.00168538 \\
2 & Colab & 1.00 & 0.00012922 & 0.00168514 \\
3 & local & 1.00 & 0.00015688 & 0.00127435 \\
3 & Colab & 1.00 & 0.00015736 & 0.00127482 \\
4 & local & 1.00 & 0.00018311 & 0.00173563 \\
4 & Colab & 1.00 & 0.00018358 & 0.00173560 \\
\hline
\end{tabular}
\end{center}

\subsection{Cross-Scale Evidence}

The cross-scale aggregate used 12 reports across seed-smoke, validation, and
extended scales. All 12 selected temporal consistency. The mean loss
improvement from $\alpha=0$ was 0.00019213. The max ordinary-logit delta was
0.00241095, and the max pinned-vs-repicked delta was 0.00556183.

\subsection{Promotion Gate}

The final promotion gate required both larger char-level and non-char tokenized
local/Colab evidence. It passed:

\begin{center}
\begin{tabular}{l l r r r}
\hline
Gate & Backend & Loss improvement & Logit delta & Support change \\
\hline
larger char & local & 0.00005436 & 0.00062868 & 0.65039062 \\
larger char & Colab & 0.00005412 & 0.00062868 & 0.65039062 \\
tokenized & local & 0.00011063 & 0.00174150 & 0.60546875 \\
tokenized & Colab & 0.00011015 & 0.00174144 & 0.60546875 \\
\hline
\end{tabular}
\end{center}

The promotion-gate satisfaction report had:

\begin{itemize}
  \item Decision: \texttt{satisfy\_temporal\_label\_free\_support\_stress\_promotion\_gate}
  \item Promotion gate satisfied: true.
  \item Promote to default support-stress mitigation: true.
  \item Mean temporal loss improvement from $\alpha=0$: 0.00008231.
  \item Max ordinary-logit delta: 0.00174150.
  \item Max pinned-vs-repicked delta: 0.00397229.
\end{itemize}

\textbf{Interpretation.} Temporal consistency is robust and safe, but its
absolute effect size is small. It is a good default stabilizer, not yet a route
to large performance gains.

\section{Residual Objective Search After Temporal HEP Promotion}

After temporal clipped HEP became the default support-stress mitigation, the
research loop returned to residual learning itself. The question was whether
the residual objective should remain supervised CE or move toward PC-style or
regularized objectives.

\subsection{PC-Style Objective}

The objective-discriminative PC-vs-supervised gate disabled the support-stress
preset so that learned residual values mattered. Both local and Colab results
kept supervised CE as the default:

\begin{center}
\begin{tabular}{l r r r}
\hline
Backend & Supervised best HEP loss & PC best HEP loss & PC minus supervised \\
\hline
local & 3.58667445 & 3.59590077 & 0.00922632 \\
Colab & 3.58667445 & 3.59590077 & 0.00922632 \\
\hline
\end{tabular}
\end{center}

PC improved its own logit-MSE objective, but its supervised CE HEP loss was
worse. The diagnostic report interpreted this as an objective-alignment issue,
not as a failed HEP sweep.

\subsection{Anchored PC}

A CE-anchored PC objective nearly closed the PC gap but still did not beat
supervised CE:

\begin{center}
\begin{tabular}{l r r r}
\hline
Backend & Supervised best & Anchored PC best & Anchored minus supervised \\
\hline
local & 3.58667445 & 3.58672714 & 0.00005269 \\
Colab & 3.58667445 & 3.58672714 & 0.00005269 \\
\hline
\end{tabular}
\end{center}

\textbf{Interpretation.} PC-style objectives are not obviously useless; the
anchored variant almost matches supervised CE. But under the current gate they
do not justify continued validation.

\subsection{CE-Adjacent Objective Variants}

Several CE-adjacent objectives were tested under matching local and Colab
artifact gates:

\begin{center}
\begin{tabular}{l r l}
\hline
Variant & Mean variant minus supervised best HEP loss & Decision \\
\hline
confidence penalty & 0.00004196 & stop validation \\
margin penalty & 0.00021887 & stop validation \\
label smoothing & 0.00021148 & stop validation \\
anchored PC & 0.00005269 & stop PC validation \\
\hline
\end{tabular}
\end{center}

These variants sometimes improved their own training objective, but they did
not improve best temporal-clipped supervised CE HEP loss.

\subsection{Focal CE}

Focal CE looked promising in the first broader sweep. Across 12 artifact-backed
comparisons, focal CE had a mean focal-minus-supervised best HEP loss of
-0.00066582 and a much lower final residual training loss. This justified a
promotion/stop gate.

The stricter seed-2 promotion gate then failed:

\begin{center}
\begin{tabular}{l l r r r}
\hline
Gate & Backend & Supervised best & Focal best & Focal minus supervised \\
\hline
char xxlarge seed 2 & local & 3.38816094 & 3.38744736 & -0.00071359 \\
char xxlarge seed 2 & Colab & 3.38816094 & 3.38744736 & -0.00071359 \\
token larger seed 2 & local & 4.14445353 & 4.14522123 & 0.00076771 \\
token larger seed 2 & Colab & 4.14445305 & 4.14522123 & 0.00076818 \\
\hline
\end{tabular}
\end{center}

The promotion-gate satisfaction report therefore stopped focal validation and
kept supervised CE as the default. The mean focal-minus-supervised best HEP
loss over the required seed-2 gate was 0.00002718.

\subsection{Train-Time Temporal-Consistency Regularization}

Train-time temporal-consistency regularization was tested as a residual
objective. It passed artifacts and invariants, but the effect size was tiny:
the best validation-scale improvement was 0.00000978 at weight 0.2, and the
mean temporal-consistency-minus-supervised best HEP loss was -0.00000298.
The branch was stopped before spending Colab time.

\textbf{Interpretation.} The current evidence argues against chasing more
small CE-adjacent objective variants unless they are tied to a sharper
mechanistic hypothesis.

\section{Newest Result: Residual Capacity vs Support Width}

The latest diagnostic asked whether residual learning is bottlenecked by
column capacity or by sparse support. It fixed the default residual objective
to supervised CE and fixed HEP to temporal consistency, then varied only:

\begin{itemize}
  \item residual column count: 12 vs 24;
  \item support width: \texttt{top\_k = 1} vs \texttt{top\_k = 2}.
\end{itemize}

The local result is the most interesting recent finding:

\begin{center}
\begin{tabular}{l r r r r}
\hline
Variant & Columns & \texttt{top\_k} & Best HEP loss & Support change \\
\hline
baseline & 12 & 1 & 3.58667445 & 0.00000000 \\
capacity & 24 & 1 & 3.58667445 & 0.00000000 \\
support width & 12 & 2 & 3.50149989 & 0.37890625 \\
capacity plus support & 24 & 2 & 3.51517272 & 0.29687500 \\
\hline
\end{tabular}
\end{center}

The decision report selected \texttt{run\_colab\_residual\_capacity\_support\_diagnostic}.
Widened support improved over the baseline best HEP CE loss by 0.08517456,
accepted $\alpha=1.0$, and stayed inside the logit-delta budget. The matching
Colab validation is now the next bounded step.

\textbf{Interpretation.} This is a much larger effect than the temporal-HEP
and CE-objective tweaks. It suggests that support sparsity may be a primary
bottleneck in the current residual-column design. Doubling columns without
widening support did nothing at the validation scale, while widening support
helped substantially.

\section{Summary of Main Experimental Lessons}

\begin{enumerate}
  \item The harness is now trustworthy enough for bounded local/Colab
  comparisons, but not a substitute for larger-scale learning evidence.
  \item HEP can help, but raw HEP must be constrained by stability budgets.
  \item Pinned support is a useful diagnostic but not a default method.
  \item Clipping solves a stability problem, not the direction-finding problem.
  \item A supervised guided oracle proves useful HEP directions exist.
  \item Temporal consistency is the best deployable label-free HEP direction
  tested so far, but its effect size is small.
  \item PC-style residual learning is not currently aligned with supervised CE
  improvement. A CE anchor nearly closes the gap but still does not win.
  \item Confidence penalty, margin penalty, label smoothing, and train-time
  temporal consistency are not promising under the current gates.
  \item Focal CE is mildly promising in some scales but inconsistent across
  tokenized seed-2 evidence.
  \item The strongest current route to larger improvement is not another loss
  tweak, but changing support structure. The \texttt{top\_k = 2} diagnostic is
  the largest recent positive signal.
\end{enumerate}

\section{Current Challenges}

\subsection{Effect Size}

Many statistically clean results are numerically tiny. Temporal clipped HEP is
consistent, but its loss improvements are around $10^{-4}$ in the promotion
gate. This is enough to justify a default safety/stability change, but not
enough to be an exciting learning improvement by itself.

\subsection{Objective Alignment}

PC-style losses can improve their own objective without improving supervised
CE. Anchored PC nearly matches supervised CE, suggesting the issue is not
purely optimization failure; it is probably objective alignment or a mismatch
between the residual representation and the downstream evaluation.

\subsection{Support Dynamics}

Support repicking is not a nuisance variable; it seems central. Pinned support
can be harmful, clipped repicking can be safe, and widened support can sharply
improve validation loss. The residual layer may need richer support dynamics
before objective tweaks matter.

\subsection{Colab Bridge Fragility}

The real-Chrome Colab bridge works well enough and has produced repeated
matching evidence, but it is still a brittle automation surface. The artifact
gate and fail-closed extraction are important; reviewers should treat any
single unverified Colab output as weaker than a checked artifact bundle.

\subsection{Scale}

The experiments are still small relative to the hoped-for conceptual target.
They are useful mechanistic probes, not yet evidence of robust neural learning
advantages at serious scale.

\section{Possible Routes Toward Larger Improvements}

The following are suggested routes for external reviewers to critique or
extend.

\subsection{Route A: Support-Width and Adaptive Support}

The \texttt{top\_k = 2} result should be validated on Colab and then scaled.
Possible follow-ups:

\begin{itemize}
  \item Repeat \texttt{top\_k = 2} at larger char and tokenized scales.
  \item Compare \texttt{top\_k = 1, 2, 4} with equal compute budgets.
  \item Add an adaptive support rule: start sparse, widen when support-change
  or uncertainty crosses a threshold.
  \item Test soft support weights rather than hard top-k selection.
  \item Track whether widened support helps because of capacity, smoother
  gradients, reduced support discontinuity, or easier HEP settling.
\end{itemize}

\subsection{Route B: Better HEP Direction Signals}

The guided oracle says that useful directions exist. Temporal consistency is
deployable but weak. Possible alternatives:

\begin{itemize}
  \item Contrastive temporal consistency across multiple augmentations.
  \item Local predictive-coding errors at hidden-state level, not just logit
  MSE.
  \item Agreement between shallow and deep residual predictions.
  \item A learned local critic that predicts whether an HEP step will improve
  held-out CE, trained only on local training traces.
  \item Energy functions that penalize support thrashing while rewarding
  stable improvement.
\end{itemize}

\subsection{Route C: Residual Architecture Changes}

Because doubling columns alone did not help, the bottleneck may be interaction
between columns and support. Possible tests:

\begin{itemize}
  \item Multi-site residual insertion rather than one insertion site.
  \item Column interaction terms or low-rank mixtures between selected columns.
  \item Residual columns with recurrent state over tokens.
  \item Per-column specialization losses to avoid redundant columns.
  \item Shared support memory over short windows rather than independent
  support selection per example.
\end{itemize}

\subsection{Route D: Objective Gates With Mechanistic Targets}

The CE-adjacent objective search produced mostly tiny or inconsistent effects.
Future objective variants should be tied to a mechanism:

\begin{itemize}
  \item objectives that directly improve support stability;
  \item objectives that encourage complementary column specialization;
  \item objectives that improve HEP alpha robustness;
  \item objectives that make top-k support less discontinuous;
  \item objectives that align PC-style local prediction error with downstream
  CE rather than hoping the alignment emerges.
\end{itemize}

\subsection{Route E: Stronger Evaluation}

To decide whether a route is real, evaluation should gradually move beyond
small validation gates:

\begin{itemize}
  \item more seeds at the current validation scale;
  \item larger tokenized runs;
  \item held-out slices that are not used in gate construction;
  \item explicit compute-matched baselines;
  \item reporting absolute loss, relative loss, and stability metrics together.
\end{itemize}

\section{Questions for External Reviewers}

\begin{enumerate}
  \item Does the \texttt{top\_k = 2} result suggest a real support bottleneck,
  or could it be an artifact of the current validation setup?
  \item What is the cleanest experiment to distinguish ``more active columns''
  from ``less discontinuous support''?
  \item Should the project prioritize adaptive support, soft support, or
  multi-site residual insertion?
  \item Is temporal consistency too weak as a label-free HEP signal, or is it a
  good base signal that needs a better architecture?
  \item How could one construct a deployable approximation to the guided HEP
  oracle without using supervised labels at inference time?
  \item Is PC-style residual learning being evaluated fairly, or is supervised
  CE too dominant as the downstream gate?
  \item If PC should matter, what local objective would better align hidden
  prediction error with CE improvement?
  \item Would a learned local energy/critic be scientifically acceptable, or
  would it obscure the mechanism?
  \item What scale is the smallest meaningful next step beyond these probes?
  \item Which metrics should become primary: CE loss, stability, support
  change, alpha robustness, or value under a compute budget?
  \item Are there known sparse mixture-of-experts or retrieval-routing results
  that should inform the support-width experiments?
  \item What is the most plausible route to a visible, not merely tiny,
  improvement?
\end{enumerate}

\section{Recommended Immediate Next Experiments}

\begin{enumerate}
  \item Run the matching Colab residual capacity/support diagnostic for the
  local \texttt{top\_k = 2} result.
  \item If Colab confirms it, define a support-width promotion/scale gate:
  larger char, tokenized larger, and at least one seed-repeat setting.
  \item Add a diagnostic that compares hard \texttt{top\_k = 2} against soft
  support with equal expected active mass.
  \item Add support-transition metrics: how often support changes, whether
  changes are beneficial, and whether HEP changes support before or after loss
  improves.
  \item Delay further CE-adjacent objective variants until support-width and
  adaptive-support questions are better understood.
\end{enumerate}

\section{Key Artifact Pointers}

The following files are useful starting points for reviewers:

\begin{verbatim}
docs/phase0_report.md
AUTOMATION_STATUS.md
results/reports/temporal_clipped_hep_promotion_gate_satisfaction/decision_report.md
results/reports/residual_objective_gate_decision/decision_report.md
results/reports/focal_residual_objective_promotion_gate_satisfaction/decision_report.md
results/reports/residual_learning_next_direction/decision_report.md
results/reports/residual_capacity_support_diagnostic_decision/decision_report.md
results/comparisons/validation_residual_capacity_support_temporal_clipped_objective_gate/summary.json
\end{verbatim}

\section{Bottom Line}

The project has learned something useful but not yet transformative. Temporal
clipped HEP is a robust label-free stabilization mechanism, but it is not large
enough by itself. Residual objective tweaks mostly wash out. The most promising
new hypothesis is that sparse support is constraining the residual layer: when
support is widened from \texttt{top\_k = 1} to \texttt{top\_k = 2}, local
validation loss improves much more than any recent objective tweak. The next
critical question is whether this support-width result survives Colab
validation and larger/tokenized gates.

\end{document}
